% PAPER_TODO 2.3 and 4.3: threats that follow from the theory and the measured studies.  Sources in comments.
\section{Threats to Validity}
\label{sec:threats}

\paragraph{Assumptions of the guarantee.}
Theorem~\ref{thm:admission} assumes i.i.d.\ trials drawn with replacement (Assumption~\ref{ass:trials}); we checked that the
harness is stateless ($2{,}400$ re-run draws, $0$ differences) but not that real task streams are identically distributed.
The shift bound $\xi_s$ (Assumption~\ref{ass:shift}) is declared, not estimated; the guarantee excludes harmful admissions only
when $2\xi_s\le\theta$, and Proposition~\ref{prop:nofree} shows that no procedure learning from trials alone can do better.
Screening only footprint-overlapping pairs relies on Assumption~\ref{ass:noninteract}, which \emph{we have not tested}: two edits in
different functions can still interact through state no footprint records.  The execution rule is a design constraint that a
deployment must enforce.

\paragraph{Construct validity.}
``Harm'' is a difference from a reference: a hand-written oracle (Studies D--H), or the developers' fix (Studies P, Q).  A
difference from the fix is not necessarily a defect: in the manual review some patches arguably behaved better than the fix.
Classifying an input as valid because the fix returns a value overestimates contract regressions when generated inputs are
stand-in objects (three of $31$ reviewed units); the sensitivity analysis keeps the conclusion.  Hand-written oracles check the
author's intent and miss side effects outside it (\texttt{cerberus}).

\paragraph{Internal validity.}
The second rater shares a model family with the input generator, and sees outputs but not inputs.  $19/55$ re-run units did not
reproduce because the model-written generator changed; they enter as bounds.  The model-written generator sees the fix's diff, so
inputs concentrate near the changed region.  Four Study Q instances failed for a tool limitation (patches that create files) and
count as unusable.  The Study Q analysis script defined the inflation ratio differently from the protocol; we report both.

\paragraph{Simulation.}
The library-level results (Section~\ref{sec:eval-library}) come from synthetic libraries whose interaction model satisfies
Assumption~\ref{ass:noninteract} by construction; they show what the design guarantees and costs under that model, not how often real
skills interact.  We did not measure footprint overlap between our two real skill families: each study bound only its own family, and
measuring overlap requires binding both families on every target, which is a new study.

\paragraph{External validity.}
The skill studies cover two hand-written rewrite families on $14$ targets; we make no claim about learned skill libraries at scale,
compositions in the wild, or retirement beyond a constructed demonstration.  In Study P half of the in-band cases come from repositories
that look like student assignments; Study Q covers public resolved submissions to one benchmark.
